

\vspace{0.5 cm}

\begin{example} Calcular $\displaystyle\iiint_{V} \dfrac{1}{y - x - z - 1}\, dx\, dy\, dz$ siendo $V$ el paralelepípedo de caras opuestas
paralelas en el cual tres de sus aristas coinciden con los vectores $(5,2,1),\,(2,4,2),\,(1,3,6)$.
\end{example}

\textbf{Soluci\'on:}
\bigskip

\textbf{1. Análisis de la región de integración $V$}

\medskip

El paralelepípedo $V$ tiene un vértice en el origen y sus tres aristas salientes coinciden con los
vectores $\mathbf{a} = (5,2,1)$, $\mathbf{b} = (2,4,2)$, $\mathbf{c} = (1,3,6)$.

\begin{figure}[htbp]
    \centering
    % Llamamos al archivo con el código TikZ extraído
    
    \centering
    % ESTA LÍNEA es la que crea el PDF en la carpeta /tikz
    \tikzsetnextfilename{ejercicio21_figura} 
    % Colores originales
    \definecolor{zzttqq}{rgb}{0.6,0.2,0}
    \definecolor{ududff}{rgb}{0.30196,0.30196,1}

    \begin{tikzpicture}[line cap=round, line join=round]
    \begin{axis}[
    view={110}{25}, % Ángulo de cámara para perspectiva 3D (azimut y elevación)
    axis lines=center,
    xlabel={$x$}, ylabel={$y$}, zlabel={$z$},
    xmin=-2, xmax=10,
    ymin=-2, ymax=10,
    zmin=-2, zmax=10,
    xtick={-2,2,4,6,8},
    ytick={-2,2,4,6,8},
    ztick={-2,2,4,6,8},
    ticklabel style={font=\scriptsize},
    width=10cm, height=10cm
    ]

    % 1. Definimos los vértices del paralelepípedo en 3D real
    \coordinate (P000) at (0,0,0);
    \coordinate (P100) at (5,2,1); % Vector u
    \coordinate (P010) at (2,4,2); % Vector v
    \coordinate (P001) at (1,3,6); % Vector w

    \coordinate (P110) at (7,6,3); % u + v
    \coordinate (P101) at (6,5,7); % u + w
    \coordinate (P011) at (3,7,8); % v + w
    \coordinate (P111) at (8,9,9); % u + v + w

    % 2. Dibujamos las aristas traseras (ocultas) punteadas
    \draw[very thick, dashed, color=zzttqq] (P000) -- (P010);
    \draw[very thick, dashed, color=zzttqq] (P010) -- (P110);
    \draw[very thick, dashed, color=zzttqq] (P010) -- (P011);

    % 3. Dibujamos las caras visibles (con relleno y opacidad)
    % Cara frontal-derecha
    \filldraw[fill=zzttqq, fill opacity=0.15, draw=zzttqq, very thick] 
    (P000) -- (P100) -- (P101) -- (P001) -- cycle;
    % Cara lateral-derecha
    \filldraw[fill=zzttqq, fill opacity=0.15, draw=zzttqq, very thick] 
    (P100) -- (P110) -- (P111) -- (P101) -- cycle;
    % Cara superior
    \filldraw[fill=zzttqq, fill opacity=0.15, draw=zzttqq, very thick] 
    (P001) -- (P101) -- (P111) -- (P011) -- cycle;

    % 4. Añadimos flechas a los vectores principales
    \draw[very thick, ->, >=Triangle, color=zzttqq] (P000) -- (P100);
    \draw[very thick, dashed, ->, >=Triangle, color=zzttqq] (P000) -- (P010);
    \draw[very thick, ->, >=Triangle, color=zzttqq] (P000) -- (P001);

    % 5. Dibujamos los puntos y etiquetas
    \fill[ududff] (P000) circle (2.5pt); 
    \fill[black!60] (P100) circle (2.5pt); 
    \fill[black!60] (P010) circle (2.5pt); 
    \fill[black!60] (P001) circle (2.5pt); 
    \fill[black!60] (P110) circle (2.5pt); 
    \fill[black!60] (P101) circle (2.5pt); 
    \fill[black!60] (P011) circle (2.5pt); 
    \fill[black!60] (P111) circle (2.5pt); 

    \end{axis}
    \end{tikzpicture}
    \caption{Representación de la región pedida.  \href{https://www.geogebra.org/m/bjnjeat8}{Ver en GeoGebra}}
 
    \label{fig:Ejemplo 8.2.7}
\end{figure}


En coordenadas cartesianas, la región $V$ es difícil de describir directamente, por lo que
recurrimos a un cambio de variables.

\bigskip

\textbf{2. Cambio de variables} \textcolor{red}{Explicar arriba el cambio lineal y referenciar}

\medskip

Definimos la transformación $T: [0,1]^3 \to V$ dada por:
\[
    T(u,v,w) = u\,\mathbf{a} + v\,\mathbf{b} + w\,\mathbf{c},
\]
es decir:
\[
    \begin{cases}
        x = 5u + 2v + w,\\[4pt]
        y = 2u + 4v + 3w,\\[4pt]
        z = \phantom{5}u + 2v + 6w.
    \end{cases}
\]
Al sustituir, la nueva región de integración es el cubo unitario:
\[
    V^* = [0,1] \times [0,1] \times [0,1].
\]

\bigskip

\textbf{3. Jacobiano de la transformación}

\medskip

La matriz jacobiana de $T$ es:
\[
    J_T(u,v,w) =
    \begin{pmatrix}
        \dfrac{\partial x}{\partial u} & \dfrac{\partial x}{\partial v} & \dfrac{\partial x}{\partial w}\\[10pt]
        \dfrac{\partial y}{\partial u} & \dfrac{\partial y}{\partial v} & \dfrac{\partial y}{\partial w}\\[10pt]
        \dfrac{\partial z}{\partial u} & \dfrac{\partial z}{\partial v} & \dfrac{\partial z}{\partial w}
    \end{pmatrix}
    =
    \begin{pmatrix}
        5 & 2 & 1 \\
        2 & 4 & 3 \\
        1 & 2 & 6
    \end{pmatrix}.
\]
Calculando el determinante resulta: $|J_T| = 72$.
\bigskip

\textbf{4. Transformación del integrando}

\medskip

Expresamos $y - x - z - 1$ en las nuevas variables:
\begin{align*}
    y - x - z - 1 &= (2u + 4v + 3w) - (5u + 2v + w) - (u + 2v + 6w) - 1\\
                  &= (2 - 5 - 1)\,u + (4 - 2 - 2)\,v + (3 - 1 - 6)\,w - 1\\
                  &= -4u - 4w - 1.
\end{align*}
Observamos que para $(u,v,w) \in [0,1]^3$ se tiene $-4u - 4w - 1 \in [-9,-1]$,
por lo que el denominador nunca se anula en $V^*$ y la integral está bien definida.

\bigskip

\textbf{5. Cálculo de la integral}

\medskip

Aplicando el Teorema de Cambio de Variables (Teorema 2.5.2):
\begin{align*}
    \iiint_{V} \frac{1}{y-x-z-1}\,dx\,dy\,dz
    &= \iiint_{V^*} \frac{1}{-4u - 4w - 1}\cdot|J_T|\,du\,dv\,dw\\[6pt]
    &= 72\int_0^1\int_0^1\int_0^1 \frac{1}{-4u-4w-1}\,du\,dv\,dw.
\end{align*}

Como el integrando no depende de $v$, la integral en $v$ es trivial:
\[
    = 72\int_0^1\int_0^1 \frac{1}{-4u-4w-1}\,du\,dw.
\]

\textbf{Integral en $u$:}
\begin{align*}
    \int_0^1 \frac{du}{-4u-4w-1}
    &= \left[-\frac{1}{4}\ln\lvert 4u+4w+1\rvert\right]_0^1\\
    &= -\frac{1}{4}\bigl(\ln(4w+5) - \ln(4w+1)\bigr)\\
    &= \frac{1}{4}\ln\!\left(\frac{4w+1}{4w+5}\right).
\end{align*}

\textbf{Integral en $w$:}
\begin{align*}
    \int_0^1 \frac{1}{4}\ln\!\left(\frac{4w+1}{4w+5}\right)dw
    &= \frac{1}{4}\left[\int_0^1\ln(4w+1)\,dw - \int_0^1\ln(4w+5)\,dw\right].
\end{align*}

Para cada integral usamos $\displaystyle\int \ln(4w+a)\,dw = \frac{(4w+a)\ln(4w+a) - (4w+a)}{4}$:

\[
    \int_0^1 \ln(4w+1)\,dw
    = \left[\frac{(4w+1)\bigl(\ln(4w+1)-1\bigr)}{4}\right]_0^1
    = \frac{5(\ln 5 - 1) - 1(\ln 1 - 1)}{4}
    = \frac{5\ln 5 - 4}{4}.
\]

\[
    \int_0^1 \ln(4w+5)\,dw
    = \left[\frac{(4w+5)\bigl(\ln(4w+5)-1\bigr)}{4}\right]_0^1
    = \frac{9(\ln 9 - 1) - 5(\ln 5 - 1)}{4}
    = \frac{9\ln 9 - 5\ln 5 - 4}{4}.
\]

Por lo tanto:
\begin{align*}
    \frac{1}{4}\left[\frac{5\ln 5 - 4}{4} - \frac{9\ln 9 - 5\ln 5 - 4}{4}\right]
    &= \frac{1}{16}\bigl(5\ln 5 - 4 - 9\ln 9 + 5\ln 5 + 4\bigr)\\
    &= \frac{10\ln 5 - 9\ln 9}{16}\\
    &= \frac{10\ln 5 - 18\ln 3}{16}\\
    &= \frac{5\ln 5 - 9\ln 3}{8}.
\end{align*}

Finalmente:
\[
    \iiint_{V} \frac{1}{y-x-z-1}\,dx\,dy\,dz
    = 72 \cdot \frac{5\ln 5 - 9\ln 3}{8}
    = \boxed{9(5\ln 5 - 9\ln 3)}.
\]